\documentclass[11pt,a4paper]{article}
\usepackage[T1]{fontenc}
\usepackage[utf8]{inputenc}
\usepackage{lmodern}
\usepackage[a4paper,margin=1in]{geometry}
\usepackage{amsmath,amssymb,amsthm,mathtools}
\usepackage{booktabs,array,graphicx}
\usepackage{microtype}
\usepackage[round,authoryear]{natbib}
\usepackage[colorlinks=true,linkcolor=blue!45!black,citecolor=blue!45!black,urlcolor=blue!45!black]{hyperref}
\usepackage{xcolor}
\setlength{\parindent}{1.35em}
\setlength{\parskip}{0.15em}
\setlength{\emergencystretch}{2em}
\allowdisplaybreaks[1]
\newtheorem{proposition}{Proposition}
\newtheorem{lemma}[proposition]{Lemma}
\newtheorem{corollary}[proposition]{Corollary}
\newtheorem{remark}[proposition]{Remark}
\theoremstyle{definition}\newtheorem{assumption}{Assumption}
\numberwithin{equation}{section}
\DeclareMathOperator{\SPE}{SPE}
\newcommand{\R}{\mathbb{R}}
\newcommand{\1}{\mathbf{1}}
\newcommand{\calA}{\mathcal{A}}
\newcommand{\calP}{\mathcal{P}}
\newcommand{\WI}{W_I}
\newcommand{\KV}{K_V}
\newcommand{\KE}{K_E}
\newcommand{\KW}{K_W}
\newcommand{\aR}{a_R^*}
\newcommand{\ProofSource}[1]{\textit{Proof source:} #1.}
\title{\textbf{Agenda Rights, Sunk Compatibility Investment,\\ and Public Standardization}}
\author{Anonymous manuscript}
\date{October 2026\\\small Working paper --- full draft, not a submission-ready proof audit}
\begin{document}
\maketitle
\begin{abstract}
International standardization can be decided after firms have sunk resources into partially compatible technologies. We study how proposal rights and voting requirements interact with such investments in a three-country Cournot oligopoly. Each firm either remains undeveloped or chooses an open technological design with two continuous attributes. Governments then decide whether to mandate an already available public interface. On a certified parameter region, we characterize the terminal outcomes of pure-strategy subgame-perfect equilibria for all admissible private investment profiles, including heterogeneous designs. When one or two affirmative votes suffice, public standardization is the only pure-equilibrium terminal outcome if and only if at least two governments can independently introduce the proposal. Unanimity instead sustains regional-standard equilibria. The institutional change need not raise world welfare. Three distinct cutoffs govern a nonmember government's veto, the efficiency of retrofitting after private investment, and the ex-ante welfare ranking of alternative investment paths. Their separation reflects irreversible design expenditure. The results are conditional on an open-technology game with a legally available public standard and do not cover proprietary standards, transfers, or mixed-strategy equilibria.
\end{abstract}
\noindent\textbf{Keywords:} compatibility standards; agenda setting; sunk investment; international coordination; oligopoly.\\
\noindent\textbf{JEL:} L15; L13; F13; D72.
\bigskip

\section{Introduction}\label{sec:introduction}

A public compatibility standard can benefit consumers and firms by placing products on a common technical interface. Yet agreement about the desirability of compatibility need not produce a common standard. Governments care about the profits of domestic producers as well as the benefits to consumers, and these interests may differ across countries. An additional institutional obstacle is easy to overlook: a proposal cannot pass a vote unless someone with the legal right to place it on the agenda chooses to do so. A standard may have enough supporters to win, but no authorized sponsor willing to introduce it.

The agenda question becomes more consequential when firms choose their technologies before the government vote. A firm investing in a partially compatible architecture cannot later recover that expenditure simply because a public interface becomes mandatory. The expectation of future intervention therefore changes what firms are willing to develop in the first place. Conversely, an observed private regional standard may give an outsider government a reason to block the public alternative even though the two founding countries favor it. To evaluate the resulting equilibrium, one must connect the initial private investment decisions to credible national decisions after \emph{every} possible technological history.

We develop a three-country oligopoly model that isolates this interaction. Each country hosts one firm and one government. Firms compete in differentiated-product Cournot markets and simultaneously decide whether to develop an open technical interface. A developer selects two continuous attributes: compatibility capacity and commonality of design. Effective private compatibility is limited by what both firms have developed, so a bilateral standard depends on the minimum investment by its two founders. Technical choices increase demand or change competitive substitution, but entail sunk setup and convex physical costs. The design space is continuous and firms choose independently; technologies need not be symmetric.

Once investments are observed, governments decide whether to mandate a single legally available public interface. The constitution specifies both which national governments can \emph{propose} a mandate and how many affirmative votes it needs. Under a successful mandate, firms decide whether to comply with the public interface or cease production. A developer in the same open technological family pays only the cost of completing its missing compatibility stock. Its original commonality investment, however, is sunk and cannot be recovered. We solve the ensuing compliance, voting, and initial investment games in pure strategies. The public mandate is enforceable, but this is not a model of compulsory licensing, patent expropriation, or governments bargaining over the design of an otherwise unavailable technology.

Our first result identifies a sharp constitutional threshold in an explicit region of economic primitives. When one or two affirmative national votes are sufficient, \emph{public standardization is the only pure-strategy subgame-perfect equilibrium terminal outcome if and only if at least two governments hold independent proposal rights}. With a sole proposal holder, the other two countries can sustain a privately developed bilateral standard, and the public standard is also supported by an equilibrium. Unanimity preserves privately developed regional equilibria regardless of how proposal rights are distributed. The threshold thus concerns the set of \emph{equilibrium outcomes}, not the uniqueness of each initial technological choice. The analysis does not rule out mixed-strategy equilibria.

Why does this institutional threshold extend beyond voting over an assumed private standard? In a three-country polity, any two authorized proposers intersect every two-government majority. The observation is elementary. But an arbitrary private technology profile need not yield two favorable votes. Some heterogeneous private standards can indeed make particular governments oppose a public interface. Our central equilibrium task is therefore to rule out private investment profiles that would both evade the public vote \emph{and} survive a firm's incentive to deviate. With two developers, excessive unilateral compatibility or commonality capital can be reduced without changing the common pair technology; at higher commonality choices, a smaller design adjustment raises private profit while leaving political nonintervention intact. With three developers, surviving private profiles must have ties in their maximal design attributes; careful analysis of the remaining asymmetric cases shows that two national governments weakly support the public mandate. Only after these private-investment restrictions are established does the institutional majority-intersection argument complete the result. Our equilibrium claim is consequently stronger than an ex-post vote calculation at a symmetric regional benchmark.

The second result concerns welfare and the timing of the mandate. We distinguish three thresholds evaluated at the privately optimal two-firm regional design. The excluded government's \emph{veto threshold} measures when it individually supports mandatory public standardization. The \emph{ex-post welfare threshold} asks whether society gains by requiring a retrofit \emph{after} the regional firms have already made their investments. The \emph{ex-ante welfare threshold} instead compares building and keeping the regional design with selecting the public interface from the start. We derive a strict ordering of these thresholds. In particular, the last two differ by exactly the cost of the private commonality investment that becomes redundant under the public interface.

That wedge has a direct interpretation. A public interface can dominate private regional development when chosen \emph{before} investment, yet it can be costly to impose after firms have sunk resources into the alternative. There is also a region in which the privately developed regional equilibrium yields higher ex-ante world welfare than the public equilibrium that becomes uniquely available when proposal rights are broadened. Thus a constitutional reform that eliminates private regional equilibria does not necessarily increase aggregate welfare. When the initial constitution supports both public and regional equilibria, an overall welfare comparison still requires an equilibrium-selection rule; we do not infer a unique realized gain from the change in the equilibrium set.

Several established lines of research bear directly on this question. \citet{KatzShapiro1985} and \citet{FarrellSaloner1988} show why private compatibility incentives and coordination institutions matter. \citet{GandalShy2001} already examines foreign-standard recognition and the formation of potentially exclusionary three-country standardization unions. \citet{Klimenko2009} shows how international governments' strategic compatibility policies can produce too much or too little compatibility, while \citet{Costinot2008} compares national treatment and mutual recognition and obtains welfare rankings that depend on the externality environment. We therefore cannot present international standardization coalitions, governmental conflicts, or a reversal of welfare rankings as new in themselves.

The closest institutional and investment ideas likewise have well-established antecedents. \citet{FarrellSimcoe2012} analyze consensus standardization with neutral intervention and endogenous ex-ante improvement incentives; their discussion explicitly considers a two-vote rule in a different governance game. \citet{Ruiz2004} examines government recognition policy alongside endogenous product characteristics. Our narrower question is whether independently held \emph{national proposal rights}, combined with a voting quorum, eliminate \emph{all privately sustainable pure-SPE terminal regional standards} when each firm can choose continuous \emph{sunk} open-technology stocks. The difference lies in the sequential equilibrium domain and the regulatory instrument, not in a claim to invent voting or product differentiation in standardization research.

The model is intentionally restrictive. Public \(O(1,0)\) is preexisting and legally accessible, private interfaces belong to the open family, and governments cannot make compensating monetary transfers or select a new public standard. The main result holds on a stated rectangle of the primitives and an additional conservative local robustness neighborhood. Those assumptions make the contribution testable but limit its interpretation. The full published proposition and proof comparisons with some of the closest earlier international-policy work also remain part of the final pre-submission novelty audit; we make no claim that the complete literature has been ruled out as a source of a more general covering theorem.

The paper proceeds as follows. Section~\ref{sec:model} defines the private technology stocks, demand system, and political sequence. Section~\ref{sec:policy} solves the public-compliance subgame and derives the government policy gains. Section~\ref{sec:equilibria} characterizes the institutional threshold over the entire continuous investment domain. Section~\ref{sec:welfare} separates national veto, ex-post retrofit efficiency, and ex-ante allocation efficiency. Section~\ref{sec:robustness} discusses the jointly perturbed primitive neighborhood, earlier literature, and the assumptions carrying the result. The appendices make the proof-critical derivations and rational certificate regions explicit.

\section{The model}\label{sec:model}

We consider three countries, indexed by \(i\in\{1,2,3\}\), each hosting one firm and one national government. Firms sell in three identical national product markets. The countries differ in the technological choices made by their firms but not in consumer preferences or the underlying demand parameters. The symmetry of national markets keeps the distributional conflict in the model tied to technology rather than country size.

\subsection{Technology choices and irreversible expenditure}\label{subsec:technology}

At date 1, firms simultaneously choose either \(N\), which entails no private development, or an open design
\begin{equation}
 O(c_i,a_i),\qquad 0\leq a_i\leq c_i\leq 1.
 \label{eq:action}
\end{equation}
The first attribute, \(c_i\), determines the potential compatibility-related demand benefit; \(a_i\) measures common design investment that changes the intensity of competition with compatible and noncompatible producers. They are real physical stocks rather than freely reversible announcements. An initial open developer spends
\begin{equation}
 I(c_i,a_i)=K+\frac{\phi}{2}(c_i^2+a_i^2),
 \qquad K>0,\quad \phi>0.
 \label{eq:initial-investment}
\end{equation}
A firm choosing \(N\) spends nothing. The investment is sunk before any regulatory vote. The feasible triangle \(a_i\leq c_i\) prevents a firm from choosing more commonality than compatibility capacity.

The institution includes a legally available public interface, \(O(1,0)\). Neither its legal existence nor its design is chosen by the governments in this paper. A private investor and the public reference use the same open technological family. Therefore, after a mandate, a firm that initially developed \(O(c_i,a_i)\) need only complete its missing compatibility capacity. Its \emph{additional} compliance expense is
\begin{equation}
 R_i^{O}=\frac{\phi}{2}(1-c_i^2).
 \label{eq:retrofit}
\end{equation}
There is no second \(K\), and the sunk cost of \(a_i\) is not reimbursed. A firm that initially chose \(N\) must instead pay \(K+\phi/2\) to become compliant. Consequently the \emph{total} expenses along the two paths to the public interface are
\begin{equation}
 C_i^M =
 \begin{cases}
  K+\phi/2, &\text{if the initial choice was }N,\\
  K+\phi(1+a_i^2)/2, &\text{if the initial choice was }O(c_i,a_i).
 \end{cases}
 \label{eq:final-cost}
\end{equation}
The difference \(\phi a_i^2/2\) is the nonrecoverable cost of an attribute that is no longer used under the public reference. This equality will matter for both initial investment deviations and welfare comparisons.

\subsection{Demand and Cournot competition}\label{subsec:demand}

Conditional on the technologies active after the policy stage, write inverse demand in each national market as
\begin{equation}
 p_i=z_i-\sum_{j=1}^3 H_{ij}q_j,\qquad
 H_{ii}=m\equiv1-v,\qquad H_{ij}=H_{ji}.
 \label{eq:inverse-demand}
\end{equation}
Here \(q_i\) is firm \(i\)'s quantity in one market; symmetric quantities are produced in all three markets. The demand matrix \(H\) is positive definite throughout the parameter domain we use. An active firm maximizes \(3p_iq_i\) before investment expense, yielding the interior Cournot condition
\begin{equation}
 2m q_i+\sum_{j\ne i}H_{ij}q_j=z_i.
 \label{eq:cournot-foc}
\end{equation}
At an interior equilibrium the firm's \emph{gross} profit summed over the three markets is \(3m q_i^2\), while consumer surplus in each national market is
\begin{equation}
 CS(q,H)=\frac{1}{2}q^{\mathsf T}Hq.
 \label{eq:cs}
\end{equation}
The three-market profit and one-market consumer-surplus terms reflect their distinct geographic incidences, not a planner weighting chosen after the fact.

We specify the technologies' effect on \(z\) and \(H\) as follows. If no two firms have developed an open design, \(z_i=1\) and \(H_{ij}=\gamma_0\) for \(i\ne j\). If exactly two firms, say 1 and 2, initially develop open designs, their effective shared stocks are
\[
 c_{12}=\min\{c_1,c_2\},\qquad a_{12}=\min\{a_1,a_2\}.
\]
They then have intercepts \(z_1=z_2=1+\rho c_{12}\), \(z_3=1\), and demand cross-effects
\begin{equation}
 H_{12}=\gamma_0+\kappa a_{12}-vc_{12},
 \qquad
 H_{13}=H_{23}=\gamma_0-\eta a_{12}.
 \label{eq:two-cross-effects}
\end{equation}
Thus increased commonality has two countervailing commercial implications: it raises the substitution between insiders through \(\kappa a_{12}\), but reduces the cross-effect with the outsider through \(\eta a_{12}\). These are reduced-form competition primitives. No claim is made that this matrix is the unique microfoundation of standards competition.

If all three firms develop open designs, each pair has \(c_{ij}=\min\{c_i,c_j\}\) and \(a_{ij}=\min\{a_i,a_j\}\). In that case
\begin{equation}
 H_{ij}=\gamma_0+\kappa a_{ij}-vc_{ij}
 \quad (i\ne j),\qquad
 z_i=1+\rho\,\max_{j\ne i}c_{ij}.
 \label{eq:three-cross-effects}
\end{equation}
There is no outsider-separation term \(-\eta a_{12}\) when all firms share the open family. This distinction is intentional: \(\eta\) measures the differentiation of a bilateral cohort \emph{relative to an outsider}. Formulas (\ref{eq:two-cross-effects})--(\ref{eq:three-cross-effects}) are part of the model, not consequences of an unmentioned global platform demand system.

If a public mandate \(M\) is adopted and all firms comply, everyone uses \(O(1,0)\). The demand parameters then reduce to
\[
 z_i=1+\rho,\qquad H_{ii}=m,\qquad
 H_{ij}=\gamma_0-v,\quad i\ne j.
\]
The resulting per-market output, firm gross profit, and consumer surplus are
\begin{align}
 q_I&=\frac{1+\rho}{2(m+\gamma_0-v)}, \label{eq:qi}\\
 \pi_I^g&=3m q_I^2,\qquad
 CS_I=\frac{3}{2}\{m+2(\gamma_0-v)\}q_I^2.
 \label{eq:public-gross}
\end{align}
It will be convenient to define the national public return \emph{net of the new compatibility stock but before setup expenditure}
\begin{equation}
 W_I\equiv\pi_I^g+CS_I-\frac{\phi}{2}.
 \label{eq:WI}
\end{equation}
A previously noninvesting country obtains \(W_I-K\) under full public compliance. A previously investing country instead obtains \(W_I-K-\phi a_i^2/2\). Both expressions include the same historical sunk investments as the continuation accounts that they are compared with.

\subsection{Government decisions and equilibrium concept}\label{subsec:timing}

At date 2, governments observe the entire vector of initial technology choices. A constitution consists of a nonempty \emph{agenda set} \(\calA\subseteq\{1,2,3\}\) and a \emph{YES quorum} \(r\in\{1,2,3\}\). Every government in \(\calA\) may independently introduce the legally available public mandate; a mandate passes if at least \(r\) governments vote YES. A government always proposes an implementable mandate if it weakly gains from doing so, and votes YES when its national welfare under the mandate weakly exceeds the status quo. These tie rules are maintained throughout. Since the possibility of enactment is common knowledge, this deterministic continuation may be represented by a sequence of individually available public motions and national votes; the legal proposal right does not confer bargaining power over the technical content.

The policy alternatives are \(L\) (no government intervention, private designs remain in place) and \(M\) (all production must use \(O(1,0)\)). After \(M\), firms simultaneously choose compliance or exit; the Cournot stage follows. Under \(L\), firms compete with their installed technologies. Each government's objective is the net profit of its domestic firm across the three markets plus consumer surplus in its own national market:
\begin{equation}
 \mathcal W_i=3m q_i^2+CS(q,H)
 -\bigl(\text{initial sunk and subsequent compliance costs paid by firm }i\bigr).
 \label{eq:government-welfare}
\end{equation}
There are no fiscal transfers across governments, no product patents or licensing royalties, and no privately negotiated recognition agreement. The equilibrium concept is a \emph{pure-strategy subgame-perfect Nash equilibrium} (SPE) over firms' continuous date-1 choices and every subsequent policy and compliance history. A final public-IS \emph{outcome} is not the same object as a unique vector of initial investments.

\subsection{The certified parameter region}\label{subsec:parameters}

Our main institutional classification holds for
\begin{equation}
 (v,\gamma_0,\kappa,\rho,\phi)
 =(.01,.14,.02,.32,.14),\qquad
 \eta\in[.10,.14],\quad K\in[.54,.56].
 \label{eq:baseline-box}
\end{equation}
The theorem is uniform in \(\eta,K\) on this rectangle and in all firms' admissible initial designs. It is \emph{not} a statement about every economically conceivable value of network effects or investment expense. Section~\ref{sec:robustness} describes a separately certified, narrow multi-parameter neighborhood.

\section{Compliance and national policy incentives}\label{sec:policy}

We begin with the continuation game after private investment. Two distinctions are central. A firm that has already invested in the open family faces a different incremental cost from a firm that has not invested. Moreover, a government's desire to replace a particular private standard is not sufficient to infer that such a standard can arise in a pure equilibrium of the \emph{preceding} investment game.

\subsection{A public mandate is implemented}\label{subsec:adoption}

Following a public mandate, an active compliant firm uses \(O(1,0)\); a noncompliant firm exits. With no other compliant producers, the firm's gross Cournot return is \(3/(4m)\). When one or two other firms comply, its operating return is, respectively,
\begin{equation}
 \pi_I^{(2)}
 =3m\left(\frac{1+\rho}{2m+\gamma_0-v}\right)^2,
 \quad
 \pi_I^{(3)}
 =3m\left(\frac{1+\rho}{2(m+\gamma_0-v)}\right)^2.
 \label{eq:adoption-profits}
\end{equation}
On (\ref{eq:baseline-box}), the isolated-adopter return is the smallest of the three. No initial history can require an \(N\) firm to spend more than \(K+\phi/2\) to comply, or a previous \(O\) firm to spend more than \(\phi/2\) in additional outlays. Since
\begin{equation}
 \frac{3}{4m}=\frac{25}{33}
 >\frac{14}{25}+\frac{7}{100},
 \label{eq:adoption-margin}
\end{equation}
adoption is strictly preferable to exit regardless of the other firms' actions.

\begin{lemma}[Mandatory compliance]\label{lem:adoption}
For every initial technology profile and every parameter in
(\ref{eq:baseline-box}), all firms' compliance with an enacted public
mandate is the unique Nash continuation of the compliance subgame.
\end{lemma}
\begin{proof}
Equation (\ref{eq:adoption-profits}) computes the operating payoff for
each possible number of other adopters. Inequality
(\ref{eq:adoption-margin}) compares their minimum with the largest
incremental installation cost faced by any firm. Sunk costs do not
change the payoff from compliance versus exit. Hence compliance is
strictly dominant for each firm, and its unique Nash profile consists
of three adopters.
\end{proof}

The induced public Cournot quantities and national welfare are therefore those in (\ref{eq:qi})--(\ref{eq:WI}), independently of the prior private designs. Let \(\mathcal W_i^L(h)\) denote country \(i\)'s payoff if intervention does not occur at the observed initial history \(h\). Define its net policy gain
\begin{equation}
 \Delta_i(h)=\mathcal W_i^M(h)-\mathcal W_i^L(h).
 \label{eq:policy-gain}
\end{equation}
Our explicit motion and tie conventions imply the public mandate passes if and only if
\begin{equation}
 \left|\{i:\Delta_i(h)\ge0\}\right|\ge r
 \quad\text{and}\quad
 \calA\cap\{i:\Delta_i(h)\ge0\}\ne\varnothing.
 \label{eq:vote-rule}
\end{equation}
This simple expression will help isolate the role of agenda power. It applies \emph{at a fixed history}; later we must ask which histories are privately supportable.

\subsection{Histories with at most one private developer}\label{subsec:zero-one}

When the initial profile has no developers or only one, there is no compatible private pair. The Cournot matrix remains at its all-\(N\) values, even if one firm has sunk an open-technology development cost. Let \(W_N^g\) be one country's initial gross firm-profit-plus-consumer-surplus value when all products have \(z_i=1\) and off-diagonal \(M_{ij}=\gamma_0\). Then
\begin{equation}
 q_N=\frac{1}{2(m+\gamma_0)},\quad
 W_N^g=3m q_N^2+
 \frac{3}{2}(m+2\gamma_0)q_N^2.
 \label{eq:WN}
\end{equation}
An initially \(N\) country's net policy gain from \(M\) is
\[
 \Delta_N=W_I-K-W_N^g.
\]
For an initial \(O(c,a)\) developer, comparing the two \emph{same-sunk-cost} histories yields
\[
 \Delta_O=W_I-W_N^g+\frac{\phi c^2}{2}
         =\Delta_N+K+\frac{\phi c^2}{2}.
\]
Both gains are strictly positive throughout (\ref{eq:baseline-box}), since \(W_I-W_N^g>K_{\max}=.56\). Thus all governments vote YES after zero or one private developer, irrespective of the constitution's nonempty agenda set and quorum.

\begin{lemma}[Unanimous demand for intervention at sparse histories]
\label{lem:sparse}
After any initial profile with zero or one privately developed open
interface, the three national governments strictly prefer the public
mandate. Hence every nonempty agenda set and every quorum \(r\le3\)
enacts \(M\).
\end{lemma}
The strict comparison is the same after any feasible \(O(c,a)\)
singleton investment: private compatibility externalities are absent
until a second open developer arrives.

\subsection{A privately optimal bilateral regional standard}\label{subsec:bilateral}

Consider firms 1 and 2 developing a common \(O(1,a)\), while firm 3 remains undeveloped. Set
\begin{equation}
 b=\gamma_0+\kappa a-v,\qquad d=\gamma_0-\eta a,
 \qquad z=1+\rho,\qquad
 D=2m(2m+b)-2d^2.
 \label{eq:bilateral-definitions}
\end{equation}
Solving (\ref{eq:cournot-foc}) for the two insiders and the outsider gives
\begin{equation}
 x(a,\eta)=\frac{2mz-d}{D},
 \qquad
 y(a,\eta)=\frac{2m+b-2dz}{D}.
 \label{eq:xy}
\end{equation}
Here \(x\) is either founder's quantity in each national market, and \(y\) is the outsider's. National consumer surplus is
\begin{equation}
 C_R(a,\eta)
 =(m+b)x^2+2dxy+\frac{m}{2}y^2.
 \label{eq:CR}
\end{equation}
The founder's total private return before the fixed setup \(K\), and the outsider's gross profit, are
\begin{equation}
 f(a,\eta)=3m x^2-\frac{\phi}{2}(1+a^2),
 \qquad
 \pi_O(a,\eta)=3m y^2.
 \label{eq:private-returns}
\end{equation}

For our parameter rectangle, an installed bilateral open cohort chooses full compatibility \(c=1\) and a unique interior commonality level
\begin{equation}
 a_R^*(\eta)\in(.25,.61),\quad
 f_a(a_R^*(\eta),\eta)=0,\quad
 f_{aa}<0,\quad \frac{da_R^*}{d\eta}>0.
 \label{eq:a-star}
\end{equation}
These claims concern the constrained bilateral design maximization, not the complete six-player game. The global one-dimensional optimum, including all design-boundary checks, is verified in Appendix~\ref{app:bilateral}.

At this observed optimal regional standard, the national policy gains take the form
\begin{equation}
 (\Delta_1,\Delta_2,\Delta_3)
 =\left(B(\eta),B(\eta),K_V(\eta)-K\right),
 \label{eq:bilateral-votes}
\end{equation}
where
\begin{align}
 B(\eta)
 &=W_I-\frac{\phi}{2}a_R^*(\eta)^2
       -f(a_R^*,\eta)-C_R(a_R^*,\eta),
       \label{eq:B}\\
 K_V(\eta)
 &=W_I-\pi_O(a_R^*,\eta)-C_R(a_R^*,\eta).
       \label{eq:KV}
\end{align}
The two founders' governments strictly prefer the public mandate, \(B(\eta)>0\), while the excluded government's gain \(K_V(\eta)-K\) is strictly negative on (\ref{eq:baseline-box}). The underlying inequalities are certified uniformly over the entire privately optimal design enclosure, not inferred from a single numerical illustration. They deliver exactly two favorable governments. Whether that majority can translate its preferences into policy depends on the proposal entitlement set.

\section{Proposal rights and equilibrium standardization}\label{sec:equilibria}

A government can favor public standardization and nevertheless fail to obtain it. With one agenda holder, an excluded country may block a proposal that both regional-standard members would approve. The significance of this veto depends, however, on whether the firms can sustain their initial choices once they anticipate the government's decision. We examine the entire date-1 investment equilibrium rather than taking a bilateral private standard as given.

\subsection{The agenda-intersection mechanism}\label{subsec:agenda-geometry}

Let \(S\subseteq\{1,2,3\}\) denote two governments that weakly approve a mandate at a fixed policy history. The presence of two favorable votes need not suffice under restricted agenda control: equation~(\ref{eq:vote-rule}) also requires \(\calA\cap S\ne\varnothing\). With three countries,
\begin{equation}
 \calA\cap S\ne\varnothing
 \quad\text{for every two-element }S
 \quad\Longleftrightarrow\quad|\calA|\ge2.
 \label{eq:hitting}
\end{equation}
This is an elementary set-intersection fact, not an independent economic theorem. It matters for standards policy only if we can prove that a winning coalition exists at every technologically supportable no-intervention equilibrium.

Table~\ref{tab:constitution} states the full result. An entry reporting RS and IS indicates the existence of an SPE with each terminal standard type; it neither counts all investment vectors nor excludes further public-IS equilibria.

\begin{table}[ht]
\centering
\caption{Pure-SPE terminal standard types under each constitution}
\label{tab:constitution}
\small
\begin{tabular}{@{}lccc@{}}
\toprule
Required YES votes & \(|\calA|=1\) & \(|\calA|=2\) & \(|\calA|=3\)\\
\midrule
\(r=1\) & RS and IS & IS only & IS only\\
\(r=2\) & RS and IS & IS only & IS only\\
\(r=3\), unanimity & RS and IS & RS and IS & RS and IS\\
\bottomrule
\end{tabular}
\vspace{2mm}
\begin{minipage}{.94\linewidth}\footnotesize
\textit{Note:} Technology is fixed at (\ref{eq:baseline-box}); results hold for all \(\eta,K\) in the stated rectangle and all legal real investments \(O(c_i,a_i)\). IS-only means all \emph{pure-SPE terminal outcomes} use the public interface, not that every initial investment vector is identical.
\end{minipage}
\end{table}

\subsection{A sharp institutional threshold}\label{subsec:main-constitutional}

\begin{proposition}[The minimum effective agenda set]
\label{prop:agenda}
Under the reduced three-country open-technology game and
(\ref{eq:baseline-box}), let \(\calA\) be any nonempty set of governments
independently entitled to propose the public mandate.
\begin{enumerate}
\item If \(r\in\{1,2\}\), every pure-strategy SPE terminates in public
      international standardization if and only if \(|\calA|\ge2\).
\item If \(r\in\{1,2\}\) and \(|\calA|=1\), a privately developed
      two-firm regional standard founded by the \emph{two non-agenda
      countries} is a pure-strategy SPE terminal outcome, alongside a
      public-IS equilibrium.
\item If \(r=3\), each of the three privately optimal regional
      two-firm coalitions can support a pure-strategy SPE with no
      public intervention, irrespective of \(\calA\), and public IS
      is also attainable in an SPE.
\end{enumerate}
The result applies to the complete continuous private initial-action
space, including unequal designs and three-firm investment. It
does not assert uniqueness of initial actions or exclude mixed SPEs.
\end{proposition}

\noindent\emph{Why the threshold is sufficient.}
Take \(r\le2\) and \(|\calA|\ge2\).
By Lemma~\ref{lem:sparse}, zero or one initial developer always
induces \(M\). A hypothetical private pure-SPE outcome with exactly
two developers cannot retain \(L\). Unequal individual design
stocks permit a strict cost-saving reduction that leaves effective
bilateral technology unchanged and cannot create a newly successful
pro-public vote. With a common design \(O(c,a)\), sufficiently low
commonality \(a\) makes both founder governments strictly support
a mandate. At higher \(a\), if the member governments instead
strictly oppose the mandate, a founder can profitably lower its
commonality while preserving nonintervention. In either case,
the candidate \(L\) investment vector cannot be a Nash equilibrium.

A three-developer candidate requires an additional argument.
Changing one firm's design now affects several pairwise links,
so a symmetric optimum cannot simply be imposed.
Nevertheless, a firm with uniquely maximal \(a_i\) has wasted
irreversible expenditure: reducing it to the next-highest
level does not change any pair minimum. Conditional on an
\(L\) outcome, an analogous reduction eliminates any uniquely
maximal \(c_i\) without provoking a new governmental majority.
The surviving design patterns must therefore tie both maxima.
They fall into two exhaustive technical-order classes.
Exact national welfare inequalities then identify at least two
governments that weakly prefer \(M\) in either class.
By (\ref{eq:hitting}), one of them holds a proposal right;
its mandate passes for \(r\le2\). Appendix~\ref{app:global}
states the detailed reduction and its boundary certificates.

\noindent\emph{Why the threshold is necessary.}
Suppose only country 1 may propose. Firms 2 and 3 can jointly
invest in \(O(1,a_R^*(\eta))\), leaving firm 1 at \(N\).
At this history the agenda holder's policy gain is
\(K_V(\eta)-K<0\), by (\ref{eq:bilateral-votes}), and therefore
it never introduces \(M\). Nonintervention is credible.
Neither founder benefits by switching to \(N\), since the
resulting singleton investment history induces \(M\) and
pays less than the regionally optimal private profit.
A change to another legal \(O(c,a)\) cannot improve the
founder's payoff while \(L\) persists, by the private
optimum and pairwise minimum technology; a change
inducing \(M\) is also unprofitable, because sunk
design expenditure is not recovered. The outside
firm cannot profit by entering with any admissible
open design: a global bound on its Cournot gross
profit falls short of the necessary setup expense.
The underlying all-real inequalities, rather than
a finite investment grid, are reproduced in
Appendix~\ref{app:bilateral}.

These deviations exhaust the date-1 unilateral alternatives.
The policy game at every other private history has a
credible pure continuation, and a passed mandate induces
strict compliance by Lemma~\ref{lem:adoption}.
Hence the complementary bilateral region is an
actual pure SPE, not just a hypothetically surviving
governmental vote.

When unanimity \(r=3\) is required, the same private
configuration survives for any founder pair: the excluded
government strictly votes NO, whatever \(\calA\).
Conversely, the initial \(NNN\) profile is an SPE
under every nonempty agenda set and every \(r\le3\).
All governments strictly support \(M\) there; a
unilateral initial developer still faces \(M\),
with the same public market return but an additional
irreversible expense \(\phi a_i^2/2\).
This completes the proof of Proposition~\ref{prop:agenda}.

\subsection{Equilibrium sets and constitutional change}\label{subsec:constitution}

The private equilibrium protected by a single outsider
agenda holder does not mean the two founding
governments would prefer regional separation.
Both would adopt the public interface if either
could introduce it. The agenda holder's own
national interest prevents that potential majority
from exercising its votes.

Granting the right to propose to any second government
removes this bottleneck. The effect goes beyond reversing
a vote at an already installed private standard:
a previously privately optimal regional investment
profile ceases to be an equilibrium of the
full sequential game. Constitutionally allocated
agenda power changes what firms choose to develop,
even when private engineering costs, markets, and
available interfaces remain unchanged.

None of this implies that a broader agenda is
world-welfare improving. With unanimity or narrow
agenda control, RS and IS equilibria coexist; a
constitutional comparison must account for the
initial sunk investments supporting each.
The relevant welfare question is not just whether
a mandate would be attractive after a regional
standard has already been installed.

\section{National vetoes, sunk costs, and world welfare}\label{sec:welfare}

Proposition~\ref{prop:agenda} compares the \emph{existence} of equilibria under different constitutional rules. It is not a normative ranking of those rules. World-welfare accounting requires us to distinguish a decision made before private investment from a decision made after a regional standard has been installed. These two comparisons differ precisely because a public mandate cannot undo a sunk technological design.

\subsection{Three cutoffs with different meanings}\label{subsec:three-cutoffs}

Continue to evaluate the privately optimal bilateral design
\[
 h_R=(O(1,\aR),O(1,\aR),N),
 \qquad \aR=\aR(\eta)>0.
\]
We study three cutoffs in the fixed initial setup cost \(K\), holding the other primitives at (\ref{eq:baseline-box}). The first is the \emph{nonmember's veto cutoff} \(K_V\) defined in (\ref{eq:KV}). The nonmember approves a public mandate at this already installed standard if and only if \(K\le K_V(\eta)\). In our rectangle \(K>K_V\), so it votes NO even though the two participating countries vote YES.

To define the two welfare cutoffs, note first that the observed regional design has aggregate private-firm returns
\(2[f(\aR,\eta)-K]+\pi_O(\aR,\eta)\), while all three national markets receive consumer surplus \(C_R(\aR,\eta)\). The world return \emph{from the start}, given the regional investment path, is
\begin{equation}
 W_{R}^{\mathrm{ante}}
 =2f(\aR,\eta)+\pi_O(\aR,\eta)
 +3C_R(\aR,\eta)-2K.
 \label{eq:wr-ante}
\end{equation}
The alternative initial path with no private investment and immediate public adoption incurs three public setup costs:
\begin{equation}
 W_{I}^{\mathrm{ante}}=3(W_I-K).
 \label{eq:wi-ante}
\end{equation}
The second world-welfare cutoff, the \emph{ex-ante allocation cutoff}, is thus
\begin{equation}
 K_W(\eta)=
 3W_I-\bigl[2f(\aR,\eta)+\pi_O(\aR,\eta)+3C_R(\aR,\eta)\bigr],
 \qquad
 W_I^{\mathrm{ante}}-W_R^{\mathrm{ante}}=K_W(\eta)-K.
 \label{eq:kw}
\end{equation}
This comparison ranks two feasible equilibrium \emph{investment paths}, not a public retrofit after one of those paths has already been followed.

Now consider the distinct policy question faced by governments when private firms 1 and 2 \emph{have already built} \(O(1,\aR)\). Their initial fixed and physical costs are sunk. Switching to the public interface gives each member country's continuing welfare
\(W_I-K-\phi(\aR)^2/2\); country 3, which built nothing, obtains \(W_I-K\). Relative to the already installed regional profile, the summed three-country gain from an ex-post mandate is
\begin{equation}
 \Delta W^{\mathrm{post}}
 =2B(\eta)+\bigl[K_V(\eta)-K\bigr]
 \equiv K_E(\eta)-K,
 \quad
 K_E(\eta):=K_V(\eta)+2B(\eta).
 \label{eq:ke}
\end{equation}
The middle cutoff \(K_E\) therefore concerns the social desirability of \emph{replacing an existing regional standard}, conditional on its sunk investments. An ex-post mandate improves global welfare if \(K<K_E\), even though the excluded government's personal vote can remain negative whenever \(K>K_V\).

\begin{proposition}[Nested political and efficiency thresholds]
\label{prop:thresholds}
For every \(\eta\in[.10,.14]\), along the unique privately
optimal bilateral design \(\aR(\eta)\), the veto,
ex-post efficiency, and ex-ante allocation cutoffs satisfy
\begin{equation}
 K_V(\eta)<K_E(\eta)<K_W(\eta),\qquad
 K_W(\eta)-K_E(\eta)=\phi\,\aR(\eta)^2>0.
 \label{eq:threshold-chain}
\end{equation}
All three cutoffs decrease strictly with \(\eta\).
For every \(K\in[.54,.56]\), each efficiency cutoff crosses \(K\)
exactly once on \(\eta\in[.10,.14]\); writing the crossings
\(\eta_E(K)\) and \(\eta_W(K)\), we have
\begin{equation}
 \eta_E(K)<\eta_W(K).
 \label{eq:eta-order}
\end{equation}
At \(K=.55\), the exact isolating intervals are
\[
 .111<\eta_E(.55)<.112,\qquad
 .118<\eta_W(.55)<.119.
\]
\end{proposition}

\begin{proof}
At the regional design the members' gain from \(M\) is
\(B(\eta)>0\), giving \(K_E-K_V=2B>0\).
Subtracting (\ref{eq:ke}) from (\ref{eq:kw}) and
using the initial and retrofit cost identities
(\ref{eq:initial-investment})--(\ref{eq:final-cost}) gives
\(K_W-K_E=\phi(\aR)^2>0\).
This establishes (\ref{eq:threshold-chain}) without
a numerical approximation.

The function \(\aR(\eta)\) increases strictly in \(\eta\)
by (\ref{eq:a-star}). Exact rational Bernstein-sign
tests show that the partial derivative of each cutoff
with respect to both \(\eta\) and \(a\) is negative on
\(\eta\in[.10,.14],\,a\in[.25,.61]\); hence
the total derivative along \(\aR(\eta)\) is negative.
The cutoffs have opposing strict signs relative to
every \(K\in[.54,.56]\) at \(\eta=.10\) and \(\eta=.14\),
so each admits exactly one root in the intervening
continuous interval. Their pointwise strict ordering
implies (\ref{eq:eta-order}). Rational brackets for
\(\aR\), validated by opposite signs of \(f_a\),
and the monotonicity of \(K_E,K_W\) in \(a\),
give the displayed two root isolations.
Appendix~\ref{app:thresholds} records the required
certificate targets and reproducible verification.
\end{proof}

\subsection{The irreversibility wedge}\label{subsec:wedge}

The distinction is economically transparent. A privately developed regional standard uses the commonality stock \(a_R^*>0\) to separate a compatible insider pair from a third producer. If the public standard is chosen at the outset, none of that commonality capital needs to be built. If the government instead imposes it after the region has developed the alternative design, the same expenditure has already been made and cannot be recovered. The public outcome then requires firms to abandon the attribute that generated the regional differentiation advantage. The opportunity cost of that abandoned design is
\(\phi(\aR)^2\) across the two regional founders. The two world-welfare comparisons differ by exactly that amount.

This yields three regions at any fixed \(K\) in our certified rectangle:
\[
\begin{array}{lll}
 \eta<\eta_E(K):
 & \text{public IS is preferred before investment}
 & \text{and an ex-post retrofit is beneficial};\\[1mm]
 \eta_E(K)<\eta<\eta_W(K):
 & \text{public IS is preferred before investment}
 & \text{but a retrofit is harmful};\\[1mm]
 \eta>\eta_W(K):
 & \text{private RS is preferred before investment}
 & \text{and a retrofit is harmful}.
\end{array}
\]
At either crossing the corresponding world-welfare difference is zero. These comparisons rank precisely the two alternatives defined above; they should not be read as a solution to an unconstrained world planner's technology-design problem.

For a concrete example, take \(\eta=.115\) and \(K=.55\).
The root of the private firm's design first-order condition
lies in the certified interval
\(.39823<\aR(.115)<.39825\).
Exact rational substitution gives
\(K_E(.115)<.55<K_W(.115)\).
Thus moving directly to the public interface would have
delivered higher world welfare than investing in and
retaining the regional standard. Yet, after firms have
made the regional investments, compelling a public
retrofit lowers realized welfare. No behavioral
assumption about equilibrium selection is needed
to establish this accounting distinction.

\subsection{Implications for institutional choice}\label{subsec:welfare-constitution}

The outsider government's opposition, the ex-post total welfare effect, and the ex-ante return from competing constitutional equilibria therefore cannot be represented by one threshold. In particular, granting a second government proposal rights when \(r\le2\) makes public IS the only pure-SPE terminal type by Proposition~\ref{prop:agenda}; it does not establish that the resulting outcome maximizes world welfare. For \(\eta=.14\) and \(K=.55\), the regional equilibrium available under narrow agenda control or unanimity has higher ex-ante world welfare than the public equilibrium. For \(\eta=.10\) and the same \(K\), the ex-ante ranking is reversed.

The narrower institution may nevertheless select the public equilibrium even when regional RS is also available. Without a separate model of equilibrium selection, the analysis cannot attach an unconditional welfare gain or loss to delegating more proposal rights. What can be said is that a broader agenda \emph{removes} a possible welfare-superior private equilibrium over part of the certified parameter rectangle. The model accordingly cautions against judging the merits of a constitutional rule solely by whether it achieves universal standardization.

\section{Robustness, mechanisms, and limits}\label{sec:robustness}

The central theorem is not a claim that the precise threshold of two agenda holders arises in every model of international standardization. Its force comes from the conjunction of two properties established in the preceding sections. First, a government cannot change the content of the public standard when exercising an agenda right: it can only bring a legally specified interface to a national vote. Second, firms understand this possibility before deciding how much irreversible compatibility capital to build. Our comparative statics hold across an explicit rectangle of the original parameters, but the public reference and political institution remain economically substantive assumptions.

\subsection{A local neighborhood of primitive parameters}\label{subsec:local}

One concern with a theory certified over a specific parameter rectangle is that its equilibrium classification could disappear under an arbitrarily small perturbation of the primitives. To examine that concern, we retain the entire original game and let all seven parameters vary in a nondegenerate product neighborhood:
\begin{equation}
\begin{gathered}
 \eta\in[.10,.14],\qquad K\in[.54,.56],\qquad
 \phi\in[.13,.30],\\
 \rho\in[.319,.321],\qquad
 \kappa\in[.019,.021],\\
 v\in[.009999,.010001],\qquad
 \gamma_0\in[.139999,.140001].
\end{gathered}
\label{eq:robust-box}
\end{equation}
For every point in this box, strict compliance, initial \(NNN\)-to-\(M\) existence, and exclusion of a privately retained pure equilibrium under \(r\le2\) and \(|\calA|\ge2\) remain valid. The first five parameters admit appreciable simultaneous perturbations relative to the last two; the intervals for \(v,\gamma_0\) are intentionally conservative. Their endpoints are not inferred economic bifurcation thresholds. They merely yield a transparent exact certificate of a genuine seven-dimensional local stability result. The constructive \emph{necessity} direction of Proposition~\ref{prop:agenda}, under which exactly one proposer can sustain the complementary bilateral RS, has not separately been re-certified over this entire seven-dimensional box.

\begin{proposition}[Local sufficiency of the agenda threshold]
\label{prop:robust}
For all primitives in (\ref{eq:robust-box}), if \(r\in\{1,2\}\)
and \(|\calA|\ge2\), every pure-strategy SPE terminal standard type
is mandatory public IS. At least one such SPE exists.
\end{proposition}

\noindent The certificate behind Proposition~\ref{prop:robust}
uses the same four corporate investment-count cases as
Proposition~\ref{prop:agenda}.
The two-developer proof requires strict signs for the
low-\(a\) national approval and high-\(a\) private derivative
polynomials. The three-developer proof requires nonnegative
policy-gain numerators in the aligned and crossed
technical-order cases. Each is represented on a unit cube
by a polynomial Bernstein basis; nonnegative coefficients
then certify the sign throughout the corresponding
continuous region. Exact-rational code verifies the
five-parameter rectangles in
\(\eta,\phi,\rho,\kappa\), together with the remaining
technology coordinates, at the endpoints of \(\phi\).
Because each sign numerator is affine in \(\phi\), endpoint
certification implies the intervening continuum.

Adding \(v\) and \(\gamma_0\) introduces a complication:
some three-developer Bernstein coefficients are exactly zero
at the boundary where all relevant private technologies
coincide with \(O(1,0)\).
Strictly positive coefficients are stable under small perturbations,
but zero coefficients require a separate argument.
The relevant boundary welfare numerators vanish identically
for \emph{all} \(v,\gamma_0\), not merely at their baseline values.
After clearing the strictly positive public-technology denominator,
every other coefficient has a rational lower bound above \(1/50\).
An explicit power-basis coefficient-norm estimate bounds
its joint \(v,\gamma_0\) perturbation below \(0.013\)
on the stated narrow intervals.
Hence none of the nonzero coefficients changes sign,
and all structural zeros remain zeros.
Appendix~\ref{app:certificates} states the algebraic objects,
normalizations, and scripts corresponding to these assertions.
The appendix is part of the proof, not an invitation to
accept a finite numerical parameter grid as sufficient.

\subsection{What carries the result}\label{subsec:mechanism}

Three features do the economic work.

\paragraph{Independent public agenda rights.}
With two authorized proposers, a two-country coalition
that gains from the mandate cannot be silenced by
a single excluded national government.
If only one government controls the agenda and that
government is outside the privately compatible pair,
its refusal blocks an otherwise successful vote.
The institutional comparison therefore changes the
set of effectively available policies, not a firm's
underlying marginal costs.

\paragraph{Continuous irreversible private design.}
An installed private technology may induce national
opposition to the public mandate. Yet it is not enough
for national opposition to exist at an arbitrary
technology profile. The firm may have paid for
compatibility or commonality that does not enter any
effective pair minimum, or may be able to trim a
costly commonality attribute without changing the
policy continuation. These private deviations sharply
restrict the histories that can be supported in a
pure initial-stage equilibrium. Without this
investment argument, the set-theoretic agenda
observation would be an incomplete equilibrium
theorem.

\paragraph{A credible public fallback.}
The public \(O(1,0)\) reference is already legally
available, and an open developer can complete its
missing compatibility stock without paying a second
fixed setup. These conditions make post-mandate
compliance strictly attractive and make the
government's policy comparison well-defined
at every private investment history. If the public
interface were patented, costly to negotiate, subject
to a design choice by governments, or usable only
through private licensing, a new continuation game
would be required. Our proof does not extend to such
systems by relabeling the current parameter values.

\subsection{Relation to existing research}\label{subsec:literature}

The literature already provides several reasons to
expect private and national standard-setting
incentives to diverge. \citet{KatzShapiro1985}
study network externalities and compatibility in
imperfect competition, and \citet{FarrellSaloner1988}
compare coordination through standard-setting
committees and markets. \citet{AuriolBenaim2000}
analyze decentralized standardization dynamics.
Thus neither network-mediated coordination problems
nor inefficient decentralization is a novel claim.

The international policy setting is particularly
close to \citet{GandalShy2001}, who study
government recognition of foreign standards and
three-country standardization unions with an excluded
nonmember. \citet{Klimenko2009} analyzes
strategic international compatibility policies
and the possibility of excessive or insufficient
compatibility, while \citet{Costinot2008} compares
mutual recognition with national-treatment rules
and finds institutional welfare rankings sensitive
to the economic environment. Our objective is
not to relabel these policy distortions as new.
The present rule variable is the set of national
governments authorized to \emph{introduce} a single
common public mandate and the quorum needed to
adopt it, after privately selected sunk
technical investments.

There is also close conceptual overlap with
\citet{FarrellSimcoe2012}, who study a
standards-setting war of attrition with neutral
intervention, screening, and ex-ante incentives
to improve proposals. Their analysis explicitly
contemplates a two-vote rule in a different
standards-governance environment.
\citet{Ruiz2004} examines governments'
foreign-standard recognition alongside
endogenous firms' product characteristics.
We therefore do not claim to be the first
to connect standard-setting rules with
innovation or to endogenize product differentiation
under government standardization policy.
The more limited proposition tested here concerns
a pure-equilibrium \emph{terminal-outcome set}
under a specific three-country national
agenda/quorum constitution and continuous
sunk open-technology stocks.

This literature comparison remains a manuscript
positioning argument, not an exhaustive
priority certification. In particular,
the full final published propositions and
proof appendices of all the closest
international-standards papers have not yet
been checked independently.
A submission decision should preserve
that qualification and complete the
remaining source audit.

\subsection{Scope of equilibrium and welfare claims}\label{subsec:limits}

The model excludes proprietary X/Y interfaces,
private recognition and licensing mechanisms,
budget-balanced intergovernmental transfers,
and endogenous public standard design.
The national tie rule selects YES and
the constitution fixes which governments
can propose. Replacing any of these may
alter continuation equilibria and invalidate
the institutional result. In addition,
the paper classifies only pure-strategy
SPE terminal standard types; it neither
enumerates all pure initial strategies
under the public equilibrium nor
rules out mixed SPEs. Finally,
the ex-ante welfare comparison ranks
two \emph{specified feasible equilibrium
investment paths}; it is not the
unconstrained first-best technology
allocation. These restrictions delimit
the intended policy interpretation.

\section{Conclusion}\label{sec:conclusion}

The right to propose a public standard and the votes needed to approve it are separate institutional instruments. In a three-country economy where firms invest in irreversible open compatibility technologies before governments decide whether to mandate a preexisting public interface, the distinction changes which private investments are supportable as pure-strategy subgame-perfect equilibria.

Our principal result is a sharp constitutional threshold on a specified region of economic primitives. If one or two affirmative votes are sufficient, granting independent proposal rights to two governments eliminates every privately retained regional-standard terminal pure equilibrium; one agenda holder cannot do so, because the excluded country may suppress a proposal that both regional founders favor. Under unanimity, regional-standard equilibria persist regardless of who can initiate a vote. The proof covers continuous heterogeneous investments by all three firms, not simply an assumed symmetric regional technology.

The same institutional reform has no unconditional efficiency ranking. The nonmember's national veto, the world-welfare effect of retrofitting after private expenditure is sunk, and the welfare ranking of alternative initial public and private investment paths have three distinct cutoffs. Their ordered difference shows why desirable public standardization at the outset can coexist with an undesirable retrofit after private investment. A rule that produces a unique public equilibrium can also remove a privately developed equilibrium with higher ex-ante world welfare.

These conclusions rest on a tightly specified public fallback, a no-transfer voting game, and a tractable compatibility technology. They do not settle standardization with licensing, patent ownership, endogenous public design, or mixed strategies. What the model isolates is a mechanism in which \emph{constitutional proposal rights, rather than the voting quota alone, govern whether firms find it worthwhile to build a regional technical alternative}. The welfare implications of that mechanism depend on when regulation intervenes relative to sunk technological commitments.

\clearpage
\bibliographystyle{apalike}
\bibliography{references}
\clearpage
\appendix
\section{Equilibrium derivations}\label{app:proofs}

This appendix provides the bridge between the game's primitives and the institutional classification in Proposition~\ref{prop:agenda}. We distinguish analytical reductions from the sign inequalities established by exact-rational Bernstein or Sturm certificates. The code and certificate correspondence are recorded in Appendix~\ref{app:certificates}. We do not rely on a finite search over firms' continuous private investment sets.

\subsection{Private bilateral Cournot design}\label{app:bilateral}

Consider two identical open founders \(O(c,a)\) and an undeveloped outsider. Write
\[
 b=\gamma_0+\kappa a-vc,\quad d=\gamma_0-\eta a,
 \quad z=1+\rho c,\quad D=2m(2m+b)-2d^2.
\]
The two symmetric insider quantities \(x\) and the outsider quantity \(y\) solve
\[
 \begin{bmatrix}2m+b&d\\ 2d&2m\end{bmatrix}
 \begin{bmatrix}x\\y\end{bmatrix}
 =
 \begin{bmatrix}z\\1\end{bmatrix}.
\]
Inverting the matrix gives
\begin{equation}
 x=\frac{2mz-d}{D},\qquad
 y=\frac{2m+b-2dz}{D},\qquad D>0.
 \label{eq:app-xy}
\end{equation}
The gross net-of-stock-but-pre-setup profit of a founder is
\[
 F(c,a;\eta)=3m x(c,a,\eta)^2-\frac{\phi}{2}(c^2+a^2).
\]
At fixed \(a\) we have the exact comparative-static identity
\[
 \frac{\partial x}{\partial c}
 =\frac{2m(\rho+vx)}{D}>0,
\quad
 \frac{\partial F}{\partial c}
 =6m x\frac{\partial x}{\partial c}-\phi c.
\]
The second derivative expression alone does not establish a global optimum. The Stage11 uniform inequality uses a conservative positive lower bound on \(x\,\partial x/\partial c\) and confirms \(\partial F/\partial c>0\) on the full feasible triangle for \(\eta\in[.10,.14]\); hence the optimal pair design has \(c=1\).

Set \(f(a,\eta)=F(1,a;\eta)\) as in (\ref{eq:private-returns}). The exact-rational Stage11 certificate establishes
\[
 f_{aa}(a,\eta)<0,\qquad f_{a\eta}(a,\eta)>0
 \quad \text{for all }a\in[0,1],\quad\eta\in[0,.14],
\]
and
\[
 f_a(.25,\eta)>0,\qquad
 f_a(.61,\eta)<0
 \quad\text{for all }\eta\in[.10,.14].
\]
Therefore a unique root of \(f_a=0\) lies inside \((.25,.61)\) and is the global founder-design optimum on the full interval. The implicit function theorem gives
\[
 \frac{da_R^*(\eta)}{d\eta}
 =-\frac{f_{a\eta}}{f_{aa}}>0.
\]
Sturm root counting and polynomial Bernstein signs certify these results over intervals rather than at sampled parameter points.

The government continuation at the optimal bilateral design follows from substituting \(f(a_R^*,\eta)\), \(\pi_O(a_R^*,\eta)\) and \(C_R(a_R^*,\eta)\) into the welfare definitions. Direct subtraction gives \((B,B,K_V-K)\), with each entry's sign verified independently on the full \(\eta,a_R^*\) enclosure. By permutation invariance the same signs apply to any founding pair.

\paragraph{Uniform firm deviations at the regional profile.}
We state the global inequalities used in the full initial-stage Nash argument. A founder's payoff under an arbitrary deviating open design and subsequent \(L\) is no larger than the best feasible paired-design payoff \(f(a_R^*,\eta)-K\): its individual design cannot raise either effective shared attribute above the corresponding pair minima without paying for that increase itself. If the deviation instead induces \(M\), its payoff is bounded above by
\[
 \pi_I^g-K-\frac{\phi}{2},
\]
because any initially open firm ultimately bears at least \(K+\phi/2\) in total public-interface costs. The verified inequality
\[
 \pi_I^g-\frac{\phi}{2}<f(a_R^*,\eta)
\]
makes that change unprofitable. A founder choosing \(N\) leaves a singleton open history, which leads to \(M\); the same bound applies.

A previously undeveloped outsider that enters with any legal \(O(c_3,a_3)\) after the founders' pair is installed obtains per-market Cournot quantity strictly below \(3/5\) in the private \(L\) continuation. Its gross three-market profit is therefore below \(3m(3/5)^2=1.0692\), and at \(K\ge.54\) its net profit is strictly below \(0.5292\). The original outsider's gross \(L\) payoff exceeds the certified lower bound \(9{,}114{,}361/16{,}635{,}300\), which is above \(0.5478\). The corresponding public-\(M\) deviation pays \(K+\phi/2\) and is likewise bounded below the original payoff. These calculations cover \emph{every} admissible unilateral \(N\)-to-\(O(c,a)\) choice. Together with government continuation existence, they establish the private RS initial equilibrium when its outsider agenda holder does not introduce \(M\), or when unanimity enables the outsider veto.

\subsection{No two-developer privately retained equilibrium}\label{app:two-investors}

Suppose exactly two firms initially invest, using \(O(c_1,a_1)\) and \(O(c_2,a_2)\), and assume the proposed terminal outcome is \(L\).

If \(a_1>a_2\), reducing \(a_1\) to \(a_2\) leaves the common attribute \(\min\{a_1,a_2\}\) unchanged. Every firm's Cournot output and every government's gain from \(M\) relative to \(L\) are unchanged, but founder 1 saves
\[
 \frac{\phi}{2}(a_1^2-a_2^2)>0
\]
under either policy. The same argument applies with the firms reversed. Therefore \(a_1=a_2=a\) in any privately stable \(L\) candidate. Given equal \(a\), if \(c_1>c_2\), founder 1 can reduce \(c_1\) to \(c_2\), preserving the effective pair technology and the \(L\) market but saving \(\phi(c_1^2-c_2^2)/2\). The deviator's national government's gain from moving to \(M\) falls weakly, and the other governments' gains remain unchanged; no new mandate can arise from a history at which \(L\) originally prevailed. Thus \(c_1=c_2=c\) also holds at any candidate \(L\)-SPE.

At identical \(O(c,a)\), the two founders have common Cournot output \(x=X/D\), with
\[
 X=2m(1+\rho c)-(\gamma_0-\eta a),\qquad
 D=2m(2m+\gamma_0+\kappa a-vc)-2(\gamma_0-\eta a)^2.
\]
Write \(C_R(c,a,\eta)\) for the national consumer surplus calculated from this \(x\) and its outsider quantity. The joint government gain for either founder is
\[
 \mathcal B(c,a,\eta)
 =W_I-\frac{\phi a^2}{2}
   -\left[3mx^2-\frac{\phi}{2}(c^2+a^2)\right]
   -C_R(c,a,\eta)
 =G_I-3mx^2-C_R-\frac{\phi}{2}(1-c^2),
\]
where \(G_I=W_I+\phi/2\) is \emph{gross} public national welfare. The exact low-design certificate is
\begin{equation}
 \mathcal B(c,a,\eta)>0
 \quad \forall\,\eta\in[.10,.14],\
 a\in[0,.75],\ c\in[0,1].
 \label{eq:app-low-a}
\end{equation}
It holds even on the rectangular superset of physically legal designs. Two favorable government votes then mandate \(M\) whenever \(|\calA|\ge2\) and \(r\le2\).

For \(a>.75\), if the same founder vote gain is nonnegative, this contradiction persists. If it is negative, both founder governments strictly prefer \(L\). Let
\[
 f_c(a,c,\eta)
 =3m x(c,a,\eta)^2-\frac{\phi}{2}(c^2+a^2).
\]
The separately verified high-design certificate is
\begin{equation}
 \frac{\partial f_c}{\partial a}(a,c,\eta)<0
 \quad \forall\,\eta\in[.10,.14],\
 a,c\in[.75,1].
 \label{eq:app-high-a}
\end{equation}
An arbitrarily small unilateral reduction of one founder's \(a\) raises its private profit; by continuity the originally negative member-government policy gains stay negative. If \(L\) prevailed before the reduction, the now lower commonality cannot force a mandate through a suddenly favorable founder majority; the no-intervention outcome persists. Hence no two-developer private \(L\)-SPE exists. At policy-gain equality the YES-on-tie convention enacts \(M\), so there is no missing knife-edge.

The code certificates for (\ref{eq:app-low-a})--(\ref{eq:app-high-a}) consist of exact-rational numerator and denominator signs in transformed unit-cube variables. Appendix~\ref{app:certificates} specifies the transformations.

\subsection{The heterogeneous three-developer cases}\label{app:global}

Now every firm has invested \(O(c_i,a_i)\). As a first reduction, if one firm uniquely possesses the highest \(a_i\), it can reduce that attribute to the next-highest \(a\) without changing any of the three pairwise minima. Its sunk capital strictly falls while every market and national policy comparison stays fixed. Such a profile cannot be a private pure equilibrium. After the maximal \(a\) is tied, a uniquely highest \(c_i\) can also be reduced to the next-highest \(c\). Because another holder of maximum \(a\) has \(a_j\le c_j\), this reduction remains feasible. Its private \(L\) payoff strictly increases and only its own country's desire for \(M\) can weakly decline, so an existing \(L\) cannot turn into \(M\). Thus a putative three-developer \(L\)-SPE must tie \emph{both} maxima.

Sort the \(c_i\) so \(c_1\le c_2=c_3\). There are two effective design-order cases, up to permutation. In the \emph{aligned} case the smallest \(c\) holder also has the smallest \(a\). Parametrize this whole feasible case by \(t,s,w,h\in[0,1]\):
\begin{equation}
 c_1=ts,\quad c_2=c_3=t,\quad
 a_1=c_1w,\quad
 d:=\min\{a_2,a_3\}=a_1+(c_2-a_1)h.
 \label{eq:aligned-param}
\end{equation}
Set
\[
 b_1=\gamma_0+\kappa a_1-vc_1,\quad
 b_2=\gamma_0+\kappa d-vc_2,\quad
 z_1=1+\rho c_1,\ z_2=1+\rho c_2,
\]
and
\[
 D_A=2m(2m+b_2)-2b_1^2,\quad
 Q_1=(2m+b_2)z_1-2b_1z_2,\quad
 Q_2=2mz_2-b_1z_1.
\]
The Cournot quantities are \(q_1=Q_1/D_A\) and \(q_2=q_3=Q_2/D_A\). Put
\[
 C_A=mQ_1^2+2mQ_2^2+4b_1Q_1Q_2+2b_2Q_2^2.
\]
Direct substitution into the national policy comparison gives the bridge equation
\begin{equation}
 2D_A^2\Delta_2
 =\left(2G_I-\phi[1-c_2^2]\right)D_A^2
 -6mQ_2^2-C_A
 =:\mathcal N_A(t,s,w,h).
 \label{eq:aligned-numerator}
\end{equation}
Here \(2D_A^2>0\). Exact nonnegative tensor-Bernstein coefficients for \(\mathcal N_A\) on \([0,1]^4\) show that government 2 weakly favors \(M\). Since \(q_3=q_2\) and \(c_3\ge c_2\), its government's extra retrofit saving implies \(\Delta_3\ge\Delta_2\ge0\). Therefore governments 2 and 3 weakly support \(M\).

In the \emph{crossed} case, after permutation, \(c_1\le c_2=c_3\) and \(a_1=a_2\ge a_3\). Parameterize it as
\[
 c_1=ts,\quad c_2=c_3=t,\quad
 a_1=a_2=c_1u,\quad a_3=a_1w,
 \qquad (t,s,u,w)\in[0,1]^4.
\]
Define
\[
 p=\gamma_0+\kappa a_1-vc_1,\quad
 r_c=\gamma_0+\kappa a_3-vc_1,\quad
 b=\gamma_0+\kappa a_3-vc_2.
\]
The symmetric three-by-three Cournot first-order matrix has off-diagonal terms \(p,r_c,b\) and diagonal \(2m\). Its determinant and adjugate quantities satisfy
\[
 D_C=8m^3+2p r_c b-2m(p^2+r_c^2+b^2)>0,
 \qquad q_i=\frac{Q_i}{D_C}.
\]
With \(Q_i\) obtained from the three-by-three adjugate and \(C_C\equiv D_C^2 CS(q,M)\times2\), the two relevant policy gains satisfy
\begin{equation}
 2D_C^2\Delta_i
 =\left(2G_I-\phi[1-c_i^2]\right)D_C^2
 -6mQ_i^2-C_C
 =:\mathcal N_{C,i}(t,s,u,w),\quad i=1,2.
 \label{eq:crossed-numerators}
\end{equation}
The exact coefficient certificates verify \(\mathcal N_{C,1}\ge0\) and \(\mathcal N_{C,2}\ge0\) throughout the \emph{full} unit cube. Hence governments 1 and 2 weakly support \(M\).

Every possible residual three-developer \(L\)-SPE is covered by one of (\ref{eq:aligned-param}) or the crossed parametrization. The two government supporters identified in each case intersect every two-member agenda set, and their proposal passes any \(r\le2\) vote. Together with the two-developer exclusion and Lemma~\ref{lem:sparse}, this completes the global sufficiency proof of Proposition~\ref{prop:agenda}. The inequalities are weak on some technology-identical boundaries, where the stipulated YES-on-tie rule is essential.

\subsection{Initial public equilibrium and the policy rule}\label{app:continuation}

For clarity, the basic motion/vote rule (\ref{eq:vote-rule}) can be reconstructed directly. At a fixed initial profile, all firms adopt the public interface after \(M\) by Lemma~\ref{lem:adoption}; therefore \(\Delta_i(h)\) is a well-defined real number. Under the specified tie rule, government \(i\) weakly approves iff \(\Delta_i(h)\ge0\). At least \(r\) such approvals are necessary; an entitled government with \(\Delta_i\ge0\) must also be present to introduce a successful motion. No government with \(\Delta_i<0\) chooses to introduce an implementable mandate. This proves the equivalence in (\ref{eq:vote-rule}), including boundary equality.

At \(NNN\), every government strictly wants \(M\), and all firms comply. An initial deviator to any feasible \(O(c_i,a_i)\) creates a singleton open history, at which all governments still enact \(M\). The deviator's public gross revenue is unchanged, and (\ref{eq:final-cost}) gives
\[
 \bigl[K+\tfrac{\phi}{2}(1+a_i^2)\bigr]
 -\bigl[K+\tfrac{\phi}{2}\bigr]
 =\tfrac{\phi a_i^2}{2}\ge0.
\]
Thus no deviation raises its payoff; \(NNN\) is a pure initial Nash profile, and the previously specified credible continuation strategies complete an SPE for every nonempty \(\calA\) and all \(r\le3\).

\subsection{Exact welfare bridge}\label{app:thresholds}

With \(W_I\) defined net of the installation stock but before \(K\), the regional member's two national welfare payoffs are
\[
 \mathcal W_{\rm member}^{L}=f+C_R-K,\qquad
 \mathcal W_{\rm member}^{M}
 =W_I-K-\frac{\phi}{2}(\aR)^2.
\]
The outsider has
\[
 \mathcal W_{\rm outsider}^{L}=\pi_O+C_R,\qquad
 \mathcal W_{\rm outsider}^{M}=W_I-K.
\]
Subtracting gives exactly the two member policy gains \(B\) and outsider gain \(K_V-K\) in (\ref{eq:bilateral-votes}). Their sum is \(2B+K_V-K=K_E-K\), while subtracting the two alternative \emph{initial} world-welfare paths gives
\[
 3(W_I-K)-[2f+\pi_O+3C_R-2K]
 =K_W-K.
\]
Consequently
\[
 K_W-K_E
 =[3W_I-2f-\pi_O-3C_R]
 -[K_V+2B]
 =\phi(\aR)^2.
\]
The same algebra exhibits why \emph{the baseline date} is not incidental: it changes whether the privately installed commonality stock is a required additional investment or a resource already spent.

\section{Exact certificates and reproducibility}\label{app:certificates}

The proofs in this paper combine conventional algebraic reductions with computer-assisted \emph{sign verification on continuous domains}. A passing script is not by itself a proof of the connection between demand primitives and economic conclusions. This appendix identifies the sign targets, coordinate domains, and the bridge from a certified polynomial inequality to the statements in the text.

\subsection{From Cournot quantities to polynomial signs}

For a candidate private investment history, the interior Cournot solution has rational quantities \(q_i=Q_i/D\), where \(D>0\) is the determinant of its positive-definite first-order matrix. The national policy gain for a developer with installed compatibility stock \(c_i\) is
\[
 \Delta_i=G_I-3m\frac{Q_i^2}{D^2}
 -CS(q,H)-\frac{\phi}{2}(1-c_i^2).
\]
The developer's fixed \(K\) cancels from its own \(M-L\) comparison: it is sunk under either policy. Multiplying by \(2D^2>0\) yields
\begin{equation}
 \mathcal N_i
 =\left(2G_I-\phi(1-c_i^2)\right)D^2
  -6mQ_i^2-2D^2 CS(q,H).
 \label{eq:certificate-object}
\end{equation}
This is the precise object certified in the aligned and crossed three-developer cases. A nonnegative \(\mathcal N_i\) implies \(\Delta_i\ge0\), not strict approval in all boundary histories. Our government YES-on-tie convention supplies the remaining implication.

For two identical developers, (\ref{eq:app-xy}) gives \(D\), \(X\), and the outsider's numerator explicitly. The low-\(a\) political certificate is \(D^2\mathcal B>0\); the high-\(a\) private design certificate is \(D^3(-\partial_a f_c)>0\). The determinant is positive throughout their transformed feasible rectangles. Together these certificates implement the low/high-\(a\) dichotomy in Appendix~\ref{app:two-investors}.

\subsection{Tensor Bernstein sign verification}\label{app:bernstein}

Let \(P(t_1,\ldots,t_d)\) be a rational polynomial defined on the full unit cube \([0,1]^d\) and of coordinatewise degrees \((n_1,\ldots,n_d)\). Write it in the monomial and tensor-Bernstein bases:
\[
 P(t)=\sum_{\alpha}p_{\alpha}t^{\alpha}
     =\sum_{\beta}b_{\beta}
       \prod_{\ell=1}^d
        \binom{n_\ell}{\beta_\ell}
        t_\ell^{\beta_\ell}(1-t_\ell)^{n_\ell-\beta_\ell}.
\]
The exact change of basis is
\begin{equation}
 b_{\beta}
 =\sum_{\alpha\le\beta}
 p_{\alpha}\prod_{\ell=1}^d
 \frac{\binom{\beta_\ell}{\alpha_\ell}}
      {\binom{n_\ell}{\alpha_\ell}}.
 \label{eq:bernstein-transform}
\end{equation}
Every tensor-Bernstein basis function is nonnegative on the unit cube, so \(b_\beta\ge0\) for all indices implies \(P(t)\ge0\) for every real point of that cube. Strictly positive coefficients yield strict positivity on the cube. The converse is not asserted: a nonnegative polynomial can have some negative Bernstein coefficients in a particular degree basis. Hence the method supplies sufficient \emph{sign certificates}, not an exact numerical boundary for every model regime.

For the two-developer test, transform
\[
 \eta=.10+.04t,\quad
 (a,c)=(.75u,w)\quad\text{for low }a,
\]
and
\[
 \eta=.10+.04t,\quad
 (a,c)=(.75+.25u,\ .75+.25w)
 \quad\text{for high }a.
\]
Each transformed cube contains the corresponding feasible triangle and the independent deviations needed by the proof. The exact baseline signs are provided by
\path{code/stage18_agenda_control_majority_coalition_spe.py}
and
\path{code/stage19_universal_agenda_majority_global_two_firm_rs_exclusion.py}.
The three-developer aligned and crossed parametrizations in Appendix~\ref{app:global} map all residual tied-maximum cases to \([0,1]^4\).
Their exact national policy numerators are those in
(\ref{eq:aligned-numerator}) and (\ref{eq:crossed-numerators}),
verified by \path{code/stage21_three_open_global_no_private_spe.py}.

The follow-up Stage24 certificate simultaneously transforms the parameter intervals for \(\rho\) and \(\kappa\), and checks two \(\phi\) endpoints. Each relevant numerator is affine in \(\phi\), so sign verification at both endpoints proves the whole intervening interval. The 5-dimensional private and 6-dimensional triple-developer coefficient checks total \(251{,}436\) exact rational entries over those two endpoints. These counts are a reproducibility property of the chosen certificate representation, not evidence that the economic content grows proportionally with the number of coefficients.

\subsection{Identically zero coefficients and the local seven-parameter box}

Extending the sign tests to variable \(v,\gamma_0\) requires care at the boundaries where the private interface already implements public \(O(1,0)\). For some three-developer expressions, Bernstein coefficients are exactly zero at these boundaries. It would be invalid to assume that any weak inequality must survive a generic coefficient perturbation. The Stage25 audit instead first clears the positive public-welfare denominator \(H^2=(1+\gamma_0-2v)^2\). It checks algebraically, with symbolic \(v,\gamma_0\), that all previously zero coefficient faces correspond to \emph{identically zero} national policy gains at a public-identical technical profile.

At the baseline \(v=.01,\gamma_0=.14\), every nonzero coefficient of the lifted sign polynomials exceeds the rational lower bound \(1/50\). The code then constructs coefficient one-norm upper bounds for the partial derivatives with respect to \(v,\gamma_0\), using the inequalities
\[
 \|P Q\|_1\le\|P\|_1\|Q\|_1,\qquad
 \|\partial_v(PQ)\|_1
 \le \|\partial_vP\|_1\|Q\|_1+\|P\|_1\|\partial_vQ\|_1.
\]
The sum of the derivative bounds for each sign polynomial is below \(13{,}000\) over a wider parameter rectangle. Since every monomial-to-Bernstein conversion coefficient belongs to \([0,1]\), the variation in any Bernstein coefficient when both \(v\) and \(\gamma_0\) move by at most \(10^{-6}\) is below \(0.013<1/50\). The positive coefficients remain positive and the structural zeros remain zero. This argument establishes the narrow additional \(v,\gamma_0\) intervals in (\ref{eq:robust-box}) without relying on a finite numerical grid.

Separate exact inequalities verify that all governments prefer \(M\) at zero- and one-developer histories in this box, that every firm strictly complies after \(M\), and that the Cournot solution remains strictly interior. The full private-equilibrium \emph{necessity} construction under a single agenda holder is not transported to the seven-parameter box, as stated in the main text.

\subsection{Welfare threshold roots and audit trail}

The Stage26 source defines \(K_V\), \(K_E\), and \(K_W\) using (\ref{eq:KV}), (\ref{eq:ke}), and (\ref{eq:kw}), and verifies symbolically the identities
\[
 K_E-K_V=2B,\qquad K_W-K_E=\phi a^2.
\]
The derivatives of each cutoff in \(\eta\) and \(a\) have exactly verified negative signs over the entire enclosure \(\eta\in[.10,.14]\), \(a\in[.25,.61]\), using both numerator and denominator Bernstein coefficients. The privately optimal root \(\aR(\eta)\) is bracketed between exact rationals with \(f_a\) positive at the lower bracket and negative at the upper bracket. Strict concavity guarantees exactly one root in each bracket; substituting the appropriate endpoints into the monotone cutoffs proves the two distinct \(\eta\) crossing intervals stated in Proposition~\ref{prop:thresholds}.

The exact scripts for the principal chain are:
\begin{quote}\small
\noindent\path{code/stage11_eta_k_phase_certificates.py}\\
\path{code/stage15_three_country_veto_minimal_continuous_spe.py}\\
\path{code/stage19_universal_agenda_majority_global_two_firm_rs_exclusion.py}\\
\path{code/stage21_three_open_global_no_private_spe.py}\\
\path{code/stage25_seven_parameter_norm_robustness.py}\\
\path{code/stage26_nested_policy_welfare_cutoffs.py}\\
\path{code/stage27_agenda_hitting_set_spe_classification.py}
\end{quote}
Their version-controlled source files fix the polynomial coordinates, rational parameter bounds, degree assignments, coefficient checks and assertion conditions. The source code is included to make the derivations reproducible; successful automated execution should not be confused with an independent economic validity or literature-novelty certification.

\subsection{Proof boundaries}

The tests do \emph{not} certify all real parameter values outside the stated boxes; a negative coefficient at an outside point would not by itself prove an economic equilibrium change. They do not rule out mixed SPEs, assume no unmodeled patents or licensing, and cannot establish the first occurrence of a similar theorem in the literature. In particular, if the model is extended to a government-endogenous public design or political transfers, the policy gains in (\ref{eq:policy-gain}) would no longer describe the same continuation game and a new equilibrium analysis would be required.

\end{document}
